\begin{enumerate}[label=(\alph*)]
\item Ta xét một chỏm băng dày. Nếu tại vị trí đó có mật độ dòng nhiệt trung bình tỏa ra ngoài qua bề mặt trái đất là $q = 0.06 \si{W/m^2}$. Tính độ dày $d$ của băng bị tan trong một năm. Biết nhiệt nóng chảy của đá là $L_i =3.4 \times 10^ 5\si{J/kg}$, khối lượng riêng của băng là $\rho_i = 917\si{kg/m^3}$.

\item  Một chỏm băng ở Nam Cực được một nhóm nhà khoa học khám phá trong một chuyến thám hiểm. Khu vực này thông thường là một cao nguyên rộng tuy nhiên lần này họ phát hiện một vùng trũng sâu xuống giống miệng núi lửa. Giả sử rằng, khu vực đó là một hình nón với độ sâu $h = 100\si{m}$, bán kính $r = 500\si{m}$. Băng ở khu vực đó dày $D=2\si{km}$.
Các nhà khoa học sau đó kết luận rằng có một vụ phun trào núi lửa nhỏ ở bên dưới chỏm băng. Khi đó một lượng magma đã thâm nhập vào trong chỏm băng và đông đặc lại khiến cho một lượng băng nào đó tan chảy. Giả sử rằng băng chỉ chuyển động thẳng đứng, magma nóng chảy hoàn toàn và ở $1200 ^{\circ} C$ ở thời điểm đầu. Coi rằng magma sau khi thâm nhập và đông đặc có hình nón ngay bên dưới khu vực bị trũng, thời gian magma dâng lên rất ngắn so với thời gian trao đổi nhiệt và dòng nhiệt chỉ theo phương thẳng đứng.
    \begin{enumerate}[label=(\roman*)]
    \item Trong trạng thái cân bằng thủy tĩnh, chỉ ra rằng nước ở dưới chỏm băng có hình dạng là một hình nón. (Gợi ý: Sử dụng nguyên lí công ảo: Trong trạng thái cân bằng, tổng tất cả công thực hiện đều triệt tiêu mọi di chuyển ảo bất kỳ của cơ hệ hay $\sum \Delta W=0$.)
    \item Giả sử rằng tại mọi thời điểm, phần nước do băng bị chảy là một hình nón. Cùng với các giả thiết bên trên, ta có thể coi rằng sự tan chảy băng được chia làm hai giai đoạn. Đầu tiên, nước ở trên bề mặt magma chưa trong trạng thái cân bằng áp suất vậy nên nó chảy hết ra ngoài (giả sử nước vẫn giữ ở nhiệt độ $0 ^{\circ} C$). Sau đó, trạng thái cân bằng thủy tĩnh được thiết lập. Khi đó nước sẽ chiếm chỗ ở bên trên chỗ bị magma thâm nhập thay vì bị chảy đi (Xem hình vẽ để thấy trạng thái cuối cùng). Sau khi cân bằng nhiệt, hãy xác định: Độ cao $H$ của đỉnh nón nước ở dưới chỏm băng, độ cao $h_1$ của magma đã thâm nhập, tổng khối lượng $m_{tot}$ của nước đã bị tan chảy và khối lượng $m'$ của nước đã chảy ra ngoài. Cho khối lượng riêng của nước: $\rho_w = 1000 \si{kg/m^3}$, khối lượng riêng của đá và magma: $\rho_r = 2900 \si{kg/m^3}$, nhiệt dung riêng của đá và magma: $c_r = 700\si{J/kg.K}$, nhiệt nóng chảy của đá và magma: $L_r = 4.2\times10^5 \si{J/kg}$.
    \end{enumerate}
\begin{figure}[H]
\centering   


% Pattern Info
 
\tikzset{
pattern size/.store in=\mcSize, 
pattern size = 5pt,
pattern thickness/.store in=\mcThickness, 
pattern thickness = 0.3pt,
pattern radius/.store in=\mcRadius, 
pattern radius = 1pt}
\makeatletter
\pgfutil@ifundefined{pgf@pattern@name@_fh87k5v2v}{
\pgfdeclarepatternformonly[\mcThickness,\mcSize]{_fh87k5v2v}
{\pgfqpoint{0pt}{0pt}}
{\pgfpoint{\mcSize+\mcThickness}{\mcSize+\mcThickness}}
{\pgfpoint{\mcSize}{\mcSize}}
{
\pgfsetcolor{\tikz@pattern@color}
\pgfsetlinewidth{\mcThickness}
\pgfpathmoveto{\pgfqpoint{0pt}{0pt}}
\pgfpathlineto{\pgfpoint{\mcSize+\mcThickness}{\mcSize+\mcThickness}}
\pgfusepath{stroke}
}}
\makeatother

% Pattern Info
 
\tikzset{
pattern size/.store in=\mcSize, 
pattern size = 5pt,
pattern thickness/.store in=\mcThickness, 
pattern thickness = 0.3pt,
pattern radius/.store in=\mcRadius, 
pattern radius = 1pt}
\makeatletter
\pgfutil@ifundefined{pgf@pattern@name@_dyu01ch0a lines}{
\pgfdeclarepatternformonly[\mcThickness,\mcSize]{_dyu01ch0a}
{\pgfqpoint{0pt}{0pt}}
{\pgfpoint{\mcSize+\mcThickness}{\mcSize+\mcThickness}}
{\pgfpoint{\mcSize}{\mcSize}}
{\pgfsetcolor{\tikz@pattern@color}
\pgfsetlinewidth{\mcThickness}
\pgfpathmoveto{\pgfpointorigin}
\pgfpathlineto{\pgfpoint{\mcSize}{0}}
\pgfusepath{stroke}}}
\makeatother
\tikzset{every picture/.style={line width=0.75pt}} %set default line width to 0.75pt        

\begin{tikzpicture}[x=0.6pt,y=0.6pt,yscale=-0.7,xscale=0.7]
%uncomment if require: \path (0,578); %set diagram left start at 0, and has height of 578

%Shape: Rectangle [id:dp9880802212598202] 
\draw  [draw opacity=0][pattern=_fh87k5v2v,pattern size=6pt,pattern thickness=0.75pt,pattern radius=0pt, pattern color={rgb, 255:red, 0; green, 0; blue, 0}] (108,376.17) -- (612,376.17) -- (612,429.5) -- (108,429.5) -- cycle ;
%Straight Lines [id:da12792252410705507] 
\draw    (109.33,89.5) -- (277.33,89.5) ;
%Straight Lines [id:da19267146330039953] 
\draw    (442.67,89.5) -- (610.67,89.5) ;
%Shape: Triangle [id:dp10900369941967214] 
\draw   (362.63,144.83) -- (449.33,89.5) -- (277.33,89.5) -- cycle ;
%Shape: Rectangle [id:dp4358280242799647] 
\draw  [color={rgb, 255:red, 255; green, 255; blue, 255 }  ,draw opacity=1 ][fill={rgb, 255:red, 255; green, 255; blue, 255 }  ,fill opacity=1 ] (280,63.5) -- (446.67,63.5) -- (446.67,90.83) -- (280,90.83) -- cycle ;
%Shape: Triangle [id:dp8508222990547505] 
\draw  [pattern=_dyu01ch0a,pattern size=6pt,pattern thickness=0.75pt,pattern radius=0pt, pattern color={rgb, 255:red, 155; green, 155; blue, 155}] (365.33,258.17) -- (450.67,374.83) -- (281.33,374.83) -- cycle ;
%Shape: Triangle [id:dp5252891056813492] 
\draw  [fill={rgb, 255:red, 0; green, 0; blue, 0 }  ,fill opacity=1 ] (366,321.5) -- (450.67,374.83) -- (281.33,374.83) -- cycle ;
%Straight Lines [id:da5639116081347827] 
\draw    (109.33,376.17) -- (612,376.17) ;
%Straight Lines [id:da4840308083293853] 
\draw    (266,310) -- (307.59,351.59) ;
\draw [shift={(309,353)}, rotate = 225] [color={rgb, 255:red, 0; green, 0; blue, 0 }  ][line width=0.75]    (10.93,-3.29) .. controls (6.95,-1.4) and (3.31,-0.3) .. (0,0) .. controls (3.31,0.3) and (6.95,1.4) .. (10.93,3.29)   ;
%Straight Lines [id:da9051843331576774] 
\draw    (224,257.5) -- (224,372) ;
\draw [shift={(224,374)}, rotate = 270] [color={rgb, 255:red, 0; green, 0; blue, 0 }  ][line width=0.75]    (10.93,-3.29) .. controls (6.95,-1.4) and (3.31,-0.3) .. (0,0) .. controls (3.31,0.3) and (6.95,1.4) .. (10.93,3.29)   ;
\draw [shift={(224,255.5)}, rotate = 90] [color={rgb, 255:red, 0; green, 0; blue, 0 }  ][line width=0.75]    (10.93,-4.9) .. controls (6.95,-2.3) and (3.31,-0.67) .. (0,0) .. controls (3.31,0.67) and (6.95,2.3) .. (10.93,4.9)   ;
%Straight Lines [id:da7454681273973618] 
\draw  [dash pattern={on 4.5pt off 4.5pt}]  (224,257) -- (364,257) ;
%Straight Lines [id:da41877009939167065] 
\draw  [dash pattern={on 4.5pt off 4.5pt}]  (366,321) -- (488,321) ;
%Straight Lines [id:da15762092461538757] 
\draw    (490,321.5) -- (490,372.5) ;
\draw [shift={(490,374.5)}, rotate = 270] [color={rgb, 255:red, 0; green, 0; blue, 0 }  ][line width=0.75]    (10.93,-3.29) .. controls (6.95,-1.4) and (3.31,-0.3) .. (0,0) .. controls (3.31,0.3) and (6.95,1.4) .. (10.93,3.29)   ;
\draw [shift={(490,319.5)}, rotate = 90] [color={rgb, 255:red, 0; green, 0; blue, 0 }  ][line width=0.75]    (10.93,-4.9) .. controls (6.95,-2.3) and (3.31,-0.67) .. (0,0) .. controls (3.31,0.67) and (6.95,2.3) .. (10.93,4.9)   ;
%Straight Lines [id:da8379069438967053] 
\draw  [dash pattern={on 4.5pt off 4.5pt}]  (365,146) -- (417,146) -- (487,146) ;
%Straight Lines [id:da8677698501887788] 
\draw    (488,91.5) -- (488,141.5) ;
\draw [shift={(488,143.5)}, rotate = 270] [color={rgb, 255:red, 0; green, 0; blue, 0 }  ][line width=0.75]    (10.93,-3.29) .. controls (6.95,-1.4) and (3.31,-0.3) .. (0,0) .. controls (3.31,0.3) and (6.95,1.4) .. (10.93,3.29)   ;
\draw [shift={(488,89.5)}, rotate = 90] [color={rgb, 255:red, 0; green, 0; blue, 0 }  ][line width=0.75]    (10.93,-4.9) .. controls (6.95,-2.3) and (3.31,-0.67) .. (0,0) .. controls (3.31,0.67) and (6.95,2.3) .. (10.93,4.9)   ;
%Straight Lines [id:da9598747504955043] 
\draw  [dash pattern={on 4.5pt off 4.5pt}]  (278,35.5) -- (279,88.5) ;
%Straight Lines [id:da725615022218316] 
\draw  [dash pattern={on 4.5pt off 4.5pt}]  (448,34.5) -- (448,90.5) ;
%Straight Lines [id:da735138912749191] 
\draw    (447,44) -- (282,44) ;
\draw [shift={(280,44)}, rotate = 360] [color={rgb, 255:red, 0; green, 0; blue, 0 }  ][line width=0.75]    (10.93,-3.29) .. controls (6.95,-1.4) and (3.31,-0.3) .. (0,0) .. controls (3.31,0.3) and (6.95,1.4) .. (10.93,3.29)   ;
\draw [shift={(449,44)}, rotate = 180] [color={rgb, 255:red, 0; green, 0; blue, 0 }  ][line width=0.75]    (10.93,-4.9) .. controls (6.95,-2.3) and (3.31,-0.67) .. (0,0) .. controls (3.31,0.67) and (6.95,2.3) .. (10.93,4.9)   ;

% Text Node
\draw (153.33,115.33) node   [align=left] {\begin{minipage}[lt]{28.11pt}\setlength\topsep{0pt}
{\fontfamily{lmss}\selectfont Băng}
\end{minipage}};
% Text Node
\draw (370,306) node   [align=left] {\begin{minipage}[lt]{28.33pt}\setlength\topsep{0pt}
{\fontfamily{lmss}\selectfont Nước}
\end{minipage}};
% Text Node
\draw (230,286) node [anchor=north west][inner sep=0.75pt]   [align=left] {{\fontfamily{lmss}\selectfont Magma}};
% Text Node
\draw (180,300.4) node [anchor=north west][inner sep=0.75pt]    {$H$};
% Text Node
\draw (499,337.4) node [anchor=north west][inner sep=0.75pt]    {$h_{1}$};
% Text Node
\draw (501,107.4) node [anchor=north west][inner sep=0.75pt]    {$h$};
% Text Node
\draw (355,20.4) node [anchor=north west][inner sep=0.75pt]    {$2r$};


\end{tikzpicture}
\caption{Mặt cắt ngang theo chiều thẳng đứng của một vùng trũng hình nón ở một chỏm băng.}
\label{fig:P5_01}
\end{figure}

\end{enumerate}
